\subsection{Method Comparison}

\parabf{Baselines.} We compare \app{} against two state-of-the-art LLM-based false positive detection methods: LLM4SA~\cite{llm4sa} and LLMDFA~\cite{llmdfa}. 
The detailed baseline settings are in Appendix \ref{ap:baseline}.

\parabf{Setup and Metrics. }We use three commonly used metrics in vulnerability detection tasks: accuracy, precision, and recall. Additionally, we introduce two metrics to evaluate the effectiveness in identifying false positives generated by static analyzers: False Positive Reduction Precision (FPR\_P) and False Positive Reduction Recall (FPR\_R), which specifically measure the precision and recall for negative samples (i.e., false positives) in the dataset.
In this section, for both \app{} and the baselines, we use GPT-4o-mini~\cite{gpt4o} as the backbone LLM. To minimize randomness, we set the temperature to 0, ensuring the LLM performs greedy decoding without sampling.

\parabf{Results.}
Table~\ref{table:rq1} compares \app{} with two baselines on \ourbench{} across three bug categories. 
Overall, \app{} significantly outperforms all baselines across all metrics. Furthermore, \app{} eliminates more false positives than the baselines (with a FPR\_R improvement of 41.1\% - 105.7\%). Particularly, \app{} precisely filters out 72\% to 96\% false alarms on different bugs. Meanwhile, \app{} still detected 42 real bugs (true positives) out of 45, achieving a high recall rate of 0.93 (with a 3.3\% - 38.8\% increase). 
Among the three bug types, the improvement for UFA detection is limited due to the relatively simple nature of the cases, where false positives arise from significant validity checks before use or reallocation after free, which do not involve complex context or path constraints. 

Table~\ref{table:rq1_2} presents the effectiveness of \app{} with three static analyzers in NPD detection. 
Among three analyzers, \app{} demonstrates the greatest advantage in eliminating false positives from CodeQL (with a FPR\_R improvement of 97.2\% - 129.0\%). This improvement is attributed to CodeQL's relatively comprehensive null pointer check rules, which lead to more complex false positives, such as the complex constraint cascade (\eg{} Example-1 in Section~\ref{sec:example}), and requires more extensive context analysis. Specifically, 44\% of CodeQL cases involve context analysis, whereas only 0.29\% and 0.36\% of cases from cppcheck and infer involve context analysis, respectively.

\begin{table}[htb]
    \centering
    \caption{Comparison on Different Bugs}\label{table:rq1}
    \footnotesize
    \begin{adjustbox}{width=1\linewidth}
        \begin{tabular}{c|c|c|c|c|c|c}
        \hline
        \textbf{Type} & \textbf{Tech.} & \textbf{Acc.} &
        \textbf{FPR\_P} & \textbf{FPR\_R} &
        \textbf{Pre.} & \textbf{Recall} \\ \hline

        \multirow{4}{*}{NPD}& \app{} & \textbf{0.75}  & \textbf{0.98} & \textbf{0.72} &\textbf{ 0.36} &  \textbf{0.93}\\
        & LLM4SA &  0.57 & 0.97 &  0.51& 0.24 & 0.90\\
        & LLMDFA & 0.39  & 0.88  & 0.35  & 0.13 & 0.67 \\\hline 
        
        \multirow{4}{*}{BOF}& \app{} & \textbf{0.86} & \textbf{1.0} & \textbf{0.86} & / & /  \\
        & LLM4SA & 0.26 & \textbf{1.0} & 0.26 & / & / \\
        & LLMDFA & 0.02  & \textbf{1.0} & 0.02 & / & /\\\hline 

        \multirow{4}{*}{UAF}& \app{} & \textbf{0.96}  & \textbf{1.0} & \textbf{0.96} & / & / \\
        & LLM4SA & 0.92  & \textbf{1.0} & 0.92 & / & / \\
        & LLMDFA & 0.32  & \textbf{1.0} & 0.32 & / & / \\\hline 

        \end{tabular}
    \end{adjustbox}
\end{table}


\begin{table}[htb]
    \centering
    \caption{Comparison on Different Analyzers}\label{table:rq1_2}
    \footnotesize
    \begin{adjustbox}{width=1\linewidth}
        \begin{tabular}{c|c|c|c|c|c|c}
        \hline
        \textbf{Analyzer} & \textbf{Tech.} & \textbf{Acc.} &
        \textbf{FPR\_P} & \textbf{FPR\_R} &
        \textbf{Pre.} & \textbf{Recall} \\ \hline

        \multirow{3}{*}{CodeQL}& \app{} & \textbf{0.75}  & \textbf{0.97} & \textbf{0.71} &\textbf{0.42} &\textbf{0.91}  \\
        & LLM4SA & 0.45  & 0.93 &  0.36& 0.24&0.88\\
        & LLMDFA & 0.38  & 0.85 & 0.31 & 0.17& 0.73\\\hline 
        
        \multirow{3}{*}{Infer}& \app{} & \textbf{0.68 } &\textbf{1.0}  &  \textbf{0.63}& \textbf{0.30}&\textbf{1.0}  \\
        & LLM4SA & 0.59  & \textbf{1.0} & 0.53 & 0.25&\textbf{1.0}\\
        & LLMDFA & 0.27  & 0.71 & 0.26 & 0.07 & 0.33 \\\hline 

        \multirow{3}{*}{Cppcheck}& \app{} & \textbf{0.79}  & \textbf{1.0} & \textbf{0.77} & 0.22 & \textbf{1.0} \\
        & LLM4SA & \textbf{0.79}  & 0.95 &  \textbf{0.77 }& 0.24 & 0.60 \\
        & LLMDFA & 0.43  & 0.93 & 0.43 & \textbf{0.54} & 0.50\\\hline 

        \end{tabular}
    \end{adjustbox}
\end{table}